\phantomsection\addcontentsline{toc}{chapter}{Anhang}
\chapter*{Anhang}


\section*{A1 Auszug aus dem Event-Log}
\addcontentsline{toc}{section}{A1 Auszug aus dem Event-Log}

Der folgende Auszug zeigt zwei repräsentative Fälle aus dem synthetischen Log: einen Happy-Path-Fall (800002) und einen Fall mit Liefersperre, Falschlieferung, Gutschrift und Nachlieferung (800347). Ein Fall mit Kreditsperre und Auftragsänderung ist in Tabelle 3 dargestellt. Alle Ressourcenkennungen sind synthetisch.

{\small
\begin{longtable}{@{}l l l l@{}}
\toprule
\textbf{Case-ID} & \textbf{Aktivität} & \textbf{Zeitstempel} & \textbf{Ressource} \\
\midrule\endhead
800002 & Auftrag angelegt & 22.09.2025 11:22:10 & VERT-11 \\
800002 & Auftrag freigegeben & 22.09.2025 11:25:10 & SYSTEM \\
800002 & Lieferung angelegt & 23.09.2025 06:30:29 & SYSTEM \\
800002 & Warenausgang gebucht & 24.09.2025 15:13:53 & LAG-03 \\
800002 & Faktura erstellt & 29.09.2025 03:10:49 & BATCH-VF06 \\
800002 & Zahlung verbucht & 03.11.2025 (tagesgenau) & FI-DEB-03 \\
800347 & Auftrag angelegt & 27.01.2026 10:44:40 & VERT-04 \\
800347 & Auftrag freigegeben & 27.01.2026 10:48:40 & SYSTEM \\
800347 & Liefersperre gesetzt & 27.01.2026 10:48:42 & SYSTEM \\
800347 & Lieferung angelegt & 03.02.2026 06:30:40 & SYSTEM \\
800347 & Warenausgang gebucht & 05.02.2026 17:54:07 & LAG-06 \\
800347 & Faktura erstellt & 09.02.2026 03:10:21 & BATCH-VF06 \\
800347 & Gutschrift erstellt & 23.02.2026 12:39:05 & VERT-04 \\
800347 & Lieferung angelegt & 24.02.2026 06:30:00 & SYSTEM \\
800347 & Warenausgang gebucht & 25.02.2026 17:09:45 & LAG-01 \\
800347 & Faktura erstellt & 02.03.2026 03:10:19 & BATCH-VF06 \\
800347 & Zahlung verbucht & 30.03.2026 (tagesgenau) & FI-DEB-01 \\
\bottomrule
\end{longtable}}


\section*{A2 Parameter des Simulationsmodells (Reproduzierbarkeit)}
\addcontentsline{toc}{section}{A2 Parameter des Simulationsmodells (Reproduzierbarkeit)}

Der Log wurde mit einem Python-Skript erzeugt (Zufallsgenerator-Seed 42, PM4Py-Version 2.7, Simulationszeitraum 01.07.2025 bis 30.06.2026, 12.000 Fälle, 84.121 Ereignisse, 11 Aktivitätstypen, Export als CSV und XES). Liegezeiten werden als Lognormalverteilungen über Arbeitstage modelliert (Wochenenden übersprungen), die Fakturierung erfolgt im wöchentlichen Sammellauf (montags 03:10 Uhr), das Zahlungsdatum ist tagesgenau modelliert (analog BSEG-AUGDT).

\begin{table}[!htb]
\centering
\small
\begin{tabularx}{\textwidth}{@{}Y l l Y@{}}
\toprule
\textbf{Muster} & \textbf{Wahrscheinlichkeit} & \textbf{Verzögerung (Median)} & \textbf{gemessen im Log} \\
\midrule
K1 Auftragsänderung (davon mehrfach) & 0,22 (0,31) & 0,6 AT je Änderung & 2.725 Fälle (22,7 \%); 814 mehrfach \\
K2 Kreditsperre + manuelle Freigabe & 0,34 & 1,8 AT & 4.086 Fälle (34,1 \%); Median 1,9 AT \\
K3 Liefersperre & 0,13 & 2,9 AT & 1.470 Fälle (12,3 \%); Median 3,9 AT bis Lieferbeleg \\
K4 Wochen-Batch-Fakturierung & alle Fälle & Sammellauf montags & Median 2,6 AT WA $\rightarrow$ Faktura \\
K5 Fakturasperre (manuelle Preise) & 0,077 & 1,5 AT & 944 Fälle (7,9 \%) \\
K6 Gutschrift /\allowbreak{} Nachlieferung & 0,045 /\allowbreak{} 0,04 & 6,0 /\allowbreak{} 1,5 AT & 532 /\allowbreak{} 471 Fälle \\
Zahlungsziel-Verhalten & -- & 38 Tage (lognormal) & Median 38 Tage Faktura $\rightarrow$ Zahlung \\
\bottomrule
\end{tabularx}
\end{table}


\section*{A3 Kennzahlen des Prozessmonitorings}
\addcontentsline{toc}{section}{A3 Kennzahlen des Prozessmonitorings}

\begin{table}[!htb]
\centering
\small
\begin{tabularx}{\textwidth}{@{}Y l l Y@{}}
\toprule
\textbf{Kennzahl} & \textbf{Ist (gemessen)} & \textbf{Zielwert nach Maßnahmen} & \textbf{Wesentlicher Hebel} \\
\midrule
Anteil Fälle mit Kreditsperre & 34,1 \% & ca. 10 \% & M1 automatisierte Kreditprüfung \\
Anteil Fälle mit Liefersperre & 12,3 \% & ca. 5 \% & M2 Stammdatenqualität \\
Liegezeit Warenausgang $\rightarrow$ Faktura (Median) & 2,6 AT & ca. 0,5 AT & M3 tägliche Fakturierung \\
Fälle mit Fakturasperre & 944 (7,9 \%) & \textless{} 350 & M4 Konditionspflege, Berechtigungen \\
Manueller Klärungsaufwand p. a. (geschätzt) & ca. 1.950 h & \textless{} 700 h & M1--M5 \\
Durchlaufzeit Auftrag $\rightarrow$ Faktura (Median) & 5,8 AT & ca. 4,5 AT & M1--M4 \\
Happy-Path-Anteil (konforme Fälle) & 38,9 \% & \textgreater{} 65 \% & alle Maßnahmen \\
Anzahl Prozessvarianten & 85 & \textless{} 40 & alle Maßnahmen \\
\bottomrule
\end{tabularx}
\end{table}

